\documentclass{article}
\usepackage{amsmath, amssymb, amsthm}
\usepackage{hyperref}
\usepackage{geometry}
\geometry{margin=1in}

\title{Erd\H{o}s Problem \#387}
\date{Last updated: 2025-08-31}
\author{Source: erdosproblems.com}

\begin{document}
\maketitle

\noindent\textbf{Status:} OPEN \quad \textbf{No prize}

\noindent\textbf{Tags:} number theory, binomial coefficients

\medskip

\noindent\textbf{Problem Statement:}

Is there an absolute constant $c&#62;0$ such that, for all $1\leq k&#60; n$, the binomial coefficient $\binom{n}{k}$ has a divisor in $(cn,n]$?




Erd\H{o}s once conjectured that $\binom{n}{k}$ must always have a divisor in $(n-k,n]$, but this was disproved by Schinzel and Erd\H{o}s \cite{Sc58}. A counterexample is given by $n=99215$ and $k=15$. Schinzel conjectured (see problem B34 of \cite{Gu04}) that, for all sufficiently large $k$ which are not prime powers, there exists an $n$ such that $\binom{n}{k}$ is not divisible by any integer in $(n-k,n]$. In \cite{Er78g} Erd\H{o}s thought that this question had a negative answer. It is easy to see that $\binom{n}{k}$ always has a divisor in $[n/k,n]$. Faulkner \cite{Fa66} proved that, if $p$ is the least prime $&#62;2k$ and $n\geq p$, then $\binom{n}{k}$ has a prime divisor $\geq p$ (except $\binom{9}{2}$ and $\binom{10}{3}$). This is discussed in problems B33 and B34 of Guy's collection \cite{Gu04}, who says that Erd\H{o}s conjectured this is true for any $c&#60;1$ (if $n$ is sufficiently large).




References

[Er78g] Erd\H{o}s, P\'{a}l, On prime factors of binomial coefficients. {II} . Mat. Lapok (1978/82), 307--316.

[Fa66] Faulkner, M., On a theorem of Sylvester and Schur . J. London Math. Soc. (1966), 107--110.

[Gu04] Guy, Richard K., Unsolved problems in number theory . (2004), xviii+437.

[Sc58] Schinzel, A., Sur un probl\`eme de {P}. {E}rd\H{o}s . Colloq. Math. (1958), 198--204.


\bigskip
\noindent\small{Source: \url{https://www.erdosproblems.com/387}}
\end{document}
